% BOOK-FILE-ID: f-b1b61a0a6237
% 付録：階差型・総和型から微積分と近似計算へ

\section{和と差から微積分へ}\label{節：和と差から微積分へ}

\sectionref{節：階差型}では差を足して元の数列を求め,\,\sectionref{節：総和型}では和の差を取って項を取り出した.
差を取ることと和を取ることは,ちょうど逆向きの操作になっている.
この関係に見覚えはないだろうか.そう,微分と積分である.
隣の項との差は,1歩進んだときの変化だった.微分は,その1歩をどこまでも小さくしたときの変化である.
ここでは,差と和の世界と,微分と積分の世界を並べて見ていこう.

\subsection{和と差の対応}
以下では,最初の項を$a_0$と書くことがある.
例えば平方の列$0,1,4,9,\ldots$から隣り合う差を取ると,\,$1,3,5,\ldots$になる.
\bookterm{sequence}{数列}$a_0,a_1,a_2,\ldots$の隣り合う項の差
\[
  \Delta a_n=a_{n+1}-a_n
\]
を一\bookterm{difference}{階差}分という.
本編の階差に$\Delta$という記号と「\bookterm{difference-operator}{差分}」の名前を付けた.
たとえば$a_n=n^2$なら
\[
  \Delta a_n=2n+1,\qquad \Delta^2a_n=2.
\]
二次式,一次式,定数の順になる.\,$\Delta^2a_n=\Delta a_{n+1}-\Delta a_n$である.

\sectionref*{節：特性方程式型}の別解でも,添字を一つずらして引くことで定数項$q$を消し,
\[
  \Delta a_{n+1}=p\,\Delta a_n
\]
という等比型を作った.
これも\bookterm{difference-operator}{差分}で扱いやすい量を取り出す例である.

逆に,\bookterm{difference-operator}{差分}を足し合わせれば元の\bookterm{sequence}{数列}を求め直せる.
\[
  \begin{aligned}
    a_n-a_0
    &=(a_1-a_0)+\cdots+(a_n-a_{n-1})\\
    &=\sum_{k=0}^{n-1}\Delta a_k.
  \end{aligned}
\]
途中の項が消える.本編の\bookterm{difference}{階差}型の和は離散的な積分,総和型の$S_{n+1}-S_n=a_{n+1}$はその逆になる.図で微積分との対応を見る.
\BookFigDifference

この考え方は,\,\sectionref{節：階差等比複合型}で引く多項式$g(n)$を探したことと同じ構造をもつ.
特に$p=1$なら,\,$g(n+1)-g(n)=f(n)$を満たす式を探す問題になる.
たとえば$S_n=\sum_{k=1}^n k^2$については
\[
  S_n-S_{n-1}=n^2,\qquad S_0=0.
\]
二次式の差を逆向きに戻すので,三次式
$S_n=An^3+Bn^2+Cn+D$を候補にする.
その差は
\[
  S_n-S_{n-1}
  =3An^2+(-3A+2B)n+(A-B+C).
\]
係数比較から$A=\frac{1}{3}$, $B=\frac{1}{2}$, $C=\frac{1}{6}$となり,\,$S_0=0$から$D=0$である.
同じ初期値と差を満たすので
\[
  S_n=\frac{n(n+1)(2n+1)}6
\]
が確かめられる.
和の公式は,差を取ると目的の式になるものを探して作れる.

\paragraphbf{付加項の次数が上がる理由}\label{節：付加項の次数が上がる理由}
本編の\sectionref*{節：階差等比複合型}では,\,$p\ne1$のとき付加項と同じ次数の不動関数を引いた.
ところが,同じように見える漸化式でも次数が1つ上がることがある.例えば
\[
  a_1=0,\qquad a_{n+1}-2a_n=n2^n
\]
では,\,$a_n=2^{n-1}b_n$と置いて両辺を$2^n$で割ると$b_{n+1}-b_n=n$となり,
\[
  b_n=\frac{n(n-1)}{2},\qquad a_n=2^{n-2}n(n-1)
\]
を得る.付加項の$n$は1次式なのに,\,$b_n$は2次式になった.なぜだろうか.

差分$\Delta$で考えると理由が見える.
$\Delta$を多項式に施すと次数が1つ下がる.微分で$x^k$が$kx^{k-1}$になるのと同じである.
逆に$\Delta Q=P$となる$Q$を探すことは,微分すると$P$になる関数,つまり原始関数を探すことにあたるので,\,$Q$の次数は$P$より1つ上がるのだ.

一般に,\,$p,q\ne0$,\,$P(n)$を多項式として
\[
  a_{n+1}-pa_n=P(n)q^n
\]
を考える.\,$Q(n)q^n$の形の不動関数を探すと,代入して$q^n$で割り,
\[
  qQ(n+1)-pQ(n)=P(n)
\]
を満たす多項式$Q$を求める問題になる.
\begin{itemize}
  \item $p\ne q$なら,左辺の最高次の係数は$q-p$倍されるだけなので,\,$Q$は$P$と同じ次数で見つかる.
  \item $p=q$なら,左辺は$q\{Q(n+1)-Q(n)\}=q\,\Delta Q(n)$となる.したがって$Q$は$P$より次数が1つ高くなければならない.
\end{itemize}
付加項の指数$q$が,付加項のない式の公比$p$と一致すると,次数が1つ上がるのである.
隣接3項型の重解で$n$が現れるのも同じ仕組みによる.これは\sectionref*{節：差分で線形漸化式を解く}で扱おう.

\BookImport{f-9925c47f754b}

\BookImport{f-46762ad20d62}

\BookImport{f-d7df492a82ce}

\subsection{差分商と微分}
ここから節末までは,微分や指数・対数の計算など数学IIIの内容を使う.初めて読むときは飛ばしてもよい.
間隔$h>0$で平方を記録すると,\,$y(t)=t^2$について
\[
  \frac{y(t+h)-y(t)}h
  =\frac{(t+h)^2-t^2}h=2t+h.
\]
となり,\,$h\to0$で傾き$2t=y'(t)$に近づく.一般に$t_n=nh$での値を$a_n$とすると,単位時間あたりの変化は
\[
  \frac{a_{n+1}-a_n}{h}
\]
である.
微分可能な関数$y$について,\textbf{時刻$t$を固定}して$h$を小さくすると,\footnote{基点も動かす近似では,例えば$y'$がある閉区間で連続なら,その区間に収まる$t,t+h$について差分商と$y'(t)$の差が一様に小さくなる.ここで使う$y(t)=t^2$はこの条件を満たす.}
\[
  \frac{y(t+h)-y(t)}h\longrightarrow y'(t)
\]
となる.格子点$t_n$での傾きもこの比で近似できるが,基点も動かす場合の一様な近似には追加条件が要る.

逆に,微分方程式$y'=F(t,y)$を数値で解くには,今の点での傾きから
\[
  a_{n+1}=a_n+hF(t_n,a_n)
\]
と次の値を予測できる.
これをオイラー法という.接線の方向へ一歩ずつ進む\bookterm{recurrence}{漸化式}で曲線を近似する.

例えば$y'=-y$, $y(0)=1$の解は$e^{-t}$で,\,$0$に近づく.
オイラー法では
\[
  a_{n+1}=(1-h)a_n,\qquad a_0=1
\]
だから,\,$n\ge1$では$a_n=(1-h)^n$である.最初の値は$a_0=1$のまま保つ.
\sectionref*{節：数列の極限}で見たように,\,$0$への収束条件は$|1-h|<1$, すなわち$0<h<2$になる.
$0<h<1$では正のまま減少し,\,$h=1$なら一回で$0$になる.
$1<h<2$では符号を交互に変えながら近づくが,\,$h=2$では$1,-1,1,-1,\ldots$と繰り返して収束しない.
さらに$h>2$では絶対値が増大する.
本当の解が減衰しても,刻み幅次第で近似の列は増大する.等比型の公比は,数値計算の安定性も教える.

また,正の時刻$t$を$N$等分して$h=\frac{t}{N}$とすれば,時刻$t$における近似値は
\[
  a_N=\left(1-\frac{t}{N}\right)^N
  \longrightarrow e^{-t}.
\]
この例では,刻みを細かくした近似の収束を指数関数の\bookterm{limit}{極限}で確かめられる.一般の保証には関数や刻み幅の条件が要る.

\subsection{不動点付近の線形化}\label{節：不動点付近の線形化}

\sectionref*{節：特性方程式型}では,不動点からのずれが$p$倍され,\,$|p|<1$なら近づいた.
\bookterm{nonlinear}{非線形}の漸化式でも,微分係数と一次近似を使えば,不動点付近でずれが何倍になるかが分かる.
一項前から次が決まる\bookterm{recurrence}{漸化式}
\[
  a_{n+1}=f(a_n)
\]
を考え,\,$\,\alpha$ を\bookterm{fixed}{不動点} $f(\alpha)=\alpha$ とする.
そこからのずれを $e_n=a_n-\alpha$ と置くと
\[
  e_{n+1}=f(\alpha+e_n)-f(\alpha)
\]
である.
$f$ が $\alpha$ で微分可能なら,微分係数の定義から
\[
  \frac{f(\alpha+e)-f(\alpha)}{e}\longrightarrow f'(\alpha)
  \qquad(e\longrightarrow0)
\]
である.
つまり,ずれ $e$ が十分小さい範囲では
\[
  f(\alpha+e)-f(\alpha)\approx f'(\alpha)e
\]
となる.
不動点の近くで曲線を接線に置き換えたのである.収束も調べられる.
もし $|f'(\alpha)|<1$ なら,\,$\,|f'(\alpha)|$ より大きく $1$ より小さい数 $c$ を選べる.
微分係数の定義により,\,$\,a_n$ が $\alpha$ に十分近ければ
\[
  |e_{n+1}|
  =|f(\alpha+e_n)-f(\alpha)|
  \le c|e_n|
\]
となる.
誤差が毎回$c$倍以下に縮むので,十分近い初期値は近くにとどまり,\,$\alpha$へ収束する.
行列の\bookterm{eigenvalue}{固有値}と同じく,\,$f'(\alpha)$が不動点の近くでずれが何倍になるかを表す.
\sectionref*{節：一次分数型を行列で解く}の一次分数型で,2つの異なる\bookterm{fixed}{不動点}がある場合には
\[
  \begin{aligned}
    f'(\alpha)
    &=\frac{ps-qr}{(r\alpha+s)^2}\\
    &=\frac{\lambda_\alpha\lambda_\beta}{\lambda_\alpha^2}
    =\frac{\lambda_\beta}{\lambda_\alpha}.
  \end{aligned}
\]
行列の固有値の方程式から$ps-qr=\lambda_\alpha\lambda_\beta$である.
比の置き換えで得た\textbf{正確な公比}と,微分による\textbf{不動点の近くでの倍数}が一致した.
\bookterm{linear}{線形}代数と微分は,1回の変化で何倍になるかを測るという点で共通している.

\subsection{ニュートン法の誤差}
\sectionref{節：ニュートン法}では,\BookMathLayout{}{\\}%
$a_{n+1}=\frac12\left(a_n+\frac{2}{a_n}\right)$の$\sqrt2$への誤差がほぼ二乗された.
一般の\bookterm{newton}{ニュートン法}でも,誤差が二乗の定数倍以下になる条件を調べよう.

方程式$g(x)=0$に対する\bookterm{newton}{ニュートン法}の漸化式は
\[
  x_{n+1}=x_n-\frac{g(x_n)}{g'(x_n)}
\]
である.
$g(\alpha)=0$, $g'(\alpha)\ne0$とし,根$\alpha$の近くで$g$は二回連続微分可能とする.
本編のずれの置き換えと同じく,\,$e_n=x_n-\alpha$とおくと
\[
  e_{n+1}
  =\frac{g'(x_n)e_n-g(x_n)}{g'(x_n)}.
\]
分子は,接線による変化と実際の変化との差である.
二階微分$g''$の絶対値がこの近くで$M$以下なら
\[
  \begin{aligned}
    g'(x_n)e_n-g(x_n)
    &=\int_\alpha^{x_n}\{g'(x_n)-g'(t)\}\,dt,\\
    |g'(x_n)e_n-g(x_n)|
    &\le \frac{M}{2}|e_n|^2.
  \end{aligned}
\]
平均値の定理で$|g'(x_n)-g'(t)|\le M|x_n-t|$と評価して積分した.\,$x_n<\alpha$でも絶対値で同じ評価になる.
また,\,$g'(\alpha)\ne0$と連続性により,十分近くでは$|g'(x_n)|\ge m$となる正の定数$m$を選べる.
したがって
\[
  |e_{n+1}|\le\frac{M}{2m}|e_n|^2.
\]
十分小さな誤差は二乗の定数倍以下に縮むので,近くにとどまり急速に根へ収束する.
本編の$\sqrt2$の例で得た正確な式
\[
  e_{n+1}=\frac{e_n^2}{2a_n}
\]
は,その具体例である.

本編では$\sqrt2$を求めたが,一般の$c>0$でも同じ計算ができる.
\[
  a_1>0,\qquad a_{n+1}=\frac12\left(a_n+\frac{c}{a_n}\right)
\]
の2つの不動点$\pm\sqrt c$からのずれの比$b_n=\frac{a_n-\sqrt c}{a_n+\sqrt c}$を取ると,\,$b_{n+1}=b_n^2$となり,
\[
  a_n=\sqrt c\cdot\frac{1+b_1^{\,2^{n-1}}}{1-b_1^{\,2^{n-1}}}
\]
を得る.\,$a_1>0$より$|b_1|<1$なので$b_n\to0$,すなわち$a_n\to\sqrt c$である.
特に$a_1>\sqrt c$なら$0<b_1<1$なので$b_1=e^{-2s}\ (s>0)$と書け,
\[
  a_n=\sqrt c\cdot\frac{\cosh(2^{n-1}s)}{\sinh(2^{n-1}s)}
\]
となる.\footnote{$0<a_1<\sqrt c$のときも,相加平均と相乗平均の関係から$a_2\ge\sqrt c$となるので,\,$n\ge2$では同じ形で表せる.}
ここでも,\,$s$にあたる量が毎回2倍されている.\sectionref*{節：双曲線関数とチェビシェフ多項式}で見た$\cosh$の倍角と同じ仕組みである.

前節の\bookterm{fixed}{不動点}の考え方でも確かめられる.
漸化式の右辺の関数$N(x)=x-\frac{g(x)}{g'(x)}$について
\[
  N'(\alpha)
  =\left.\frac{g(x)g''(x)}{g'(x)^2}\right|_{x=\alpha}=0.
\]
つまり,根のところでは一次のずれが消え,次の誤差は二乗の定数倍以下になる.
二次の項まで消えると,さらに速く近づく場合がある.\footnote{例えば$g(x)=x+x^3$の根$0$では,\,$N(x)=\frac{2x^3}{1+3x^2}$となり,誤差は三次の大きさになる.\,$g(x)=x$なら一回で根に移る.}
公比が一定の等比型と違い,\bookterm{newton}{ニュートン法}では誤差が小さいほど残る割合も小さくなる.

接線の傾きを毎回微分で求めるのが難しいなら,過去二点を結ぶ割線の傾きで代用できる.
その漸化式は
\[
  x_{n+1}
  =x_n-g(x_n)\frac{x_n-x_{n-1}}{g(x_n)-g(x_{n-1})}
\]
となり,割線法と呼ばれる.
分母が0でない範囲で用いる.
直前の二項を状態にする点は前節と共通する.\bookterm{recurrence}{漸化式}は,近似の規則を作り収束を調べる道具にもなる.
